\section{Artin 代数化}\label{src:C.6}
\begin{flushright}\itshape
思考无非是看见某个引人注意的事物, 它使你看见此前未曾注意的事物,\\
后者又使你看见甚至不可见的事物.\\
\normalfont Norman Maclean, \textit{A River Through It}
\end{flushright}
Artin 代数化断言, 每个有效的万有形式形变均可 ``代数化'', 即延拓为\textit{有限型} $\Bbbk$-概形上的对象. 本节遵循 Conrad 与 de Jong 的思想 \cite{CdJ02}, 说明如何由 Artin 逼近推出 Artin 代数化.
\subsection{保持极限的预叠}\label{src:C.6.1-subsection}
推广 \S\bookref{B.5.5-subsection} 中保持极限的函子的定义, 若对每个 $\Bbbk$-代数系统 $B_\lambda$, 自然函子
\[
\operatorname{colim}\mathcal X(B_\lambda)\to\mathcal X(\operatorname{colim}B_\lambda)
\]
为范畴等价, 则称 $\Sch/\Bbbk$ 上的预叠 $\mathcal X$ \newterm{保持极限}{limit preserving} (或\textit{局部有限表现}). 当 $\mathcal X$ 为 $\Bbbk$ 上的代数叠时, 这等价于态射 $\mathcal X\to\Spec\Bbbk$ 局部有限表现; 见习题 \bookref{4.3.33}.
\begin{exercise}\label{src:C.6.1}
用 \S\bookref{B.3} 的极限方法证明, $\Hilb(X)$, $\mathcal M_g^{\mathrm{all}}$ 和 $\Coh(X)$ 均在 $\Sch/\Bbbk$ 上保持极限.
\end{exercise}
\subsection{Conrad--de Jong 逼近}\label{src:C.6.2-subsection}
Artin 逼近 (\bookref{B.5.18}) 的初始数据是完备 Noether 局部 $\Bbbk$-代数上的对象, 并假设该代数是有限生成 $\Bbbk$-代数在极大理想处的完备化. 现在我们将看到, 去掉后一个假设后, 类似的逼近结果仍成立. 思路是同时逼近完备局部环与对象.

还记得, 若 $(A,\mathfrak m)$ 为局部环, $M$ 为 $A$-模, 则 $M$ 的\textit{伴随分次模} 定义为 $\operatorname{Gr}_{\mathfrak m}(M)=\bigoplus_{n\geq0}\mathfrak m^nM/\mathfrak m^{n+1}M$; 它是分次环 $\operatorname{Gr}_{\mathfrak m}(A)$ 上的分次模.
\begin{theorem}[Conrad--de Jong 逼近]\label{src:C.6.2}
设 $\mathcal X$ 为 $\Sch/\Bbbk$ 上保持极限的预叠, $(R,\mathfrak m_R)$ 为剩余域为 $\Bbbk$ 的完备 Noether 局部 $\Bbbk$-代数, $\xi\in\mathcal X(R)$. 则对每个整数 $N\geq0$, 存在
\begin{enumerate}[(1)]
\item $\Bbbk$ 上的有限型仿射概形 $\Spec A$ 及 $\Bbbk$-点 $u\in\Spec A$;
\item 对象 $\eta\in\mathcal X(A)$;
\item 同构 $\phi_N:R/\mathfrak m_R^{N+1}\cong A/\mathfrak m_u^{N+1}$;
\item 经由 $\phi_N$ 的 $\xi|_{R/\mathfrak m_R^{N+1}}$ 与 $\eta|_{A/\mathfrak m_u^{N+1}}$ 的同构; 以及
\item 分次 $\Bbbk$-代数的同构 $\operatorname{Gr}_{\mathfrak m_R}(R)\cong\operatorname{Gr}_{\mathfrak m_u}(A)$.
\end{enumerate}
\end{theorem}
本定理的证明将同时逼近定义 $R$ 与对象 $\xi$ 的方程和关系. 断言 (1)--(4) 可容易地作为 Artin 逼近的推论得到. Conrad 与 de Jong 的巧妙洞见是, Artin 逼近还能保证条件 (5), 而且该条件足以推出 Artin 代数化中的完备局部 $\Bbbk$-代数同构. 毫不意外, 建立 (5) 最为费力.

我们需要一些预备结果, 控制 Artin--Rees 引理 (\bookref{B.5.4}) 中出现的常数.
\begin{definition}[Artin--Rees 条件]\label{src:C.6.3}
设 $(A,\mathfrak m)$ 为 Noether 局部环, $\varphi:M\to N$ 为有限 $A$-模的态射, $c\geq0$ 为整数. 若
\[
\varphi(M)\cap\mathfrak m^nN\subseteq\varphi(\mathfrak m^{n-c}M),\qquad\forall n\geq c,
\]
则称 $\varphi$ 满足 $(\mathrm{AR})_c$.
\end{definition}
Artin--Rees 引理 (\bookref{B.5.4}) 蕴含, $c\gg0$ 时 $\varphi$ 满足 $(\mathrm{AR})_c$.
\begin{lemma}\label{src:C.6.4}
设 $(A,\mathfrak m)$ 为 Noether 局部环. 设
\[
L\xrightarrow\alpha M\xrightarrow\beta N\qquad\text{和}\qquad L'\xrightarrow{\alpha'}M\xrightarrow{\beta'}N
\]
为有限 $A$-模的两个复形. 取正整数 $c$. 假设
\begin{enumerate}[(a)]
\item 第一个序列正合,
\item 两个复形模 $\mathfrak m^{c+1}$ 同构, 且
\item $\alpha$ 和 $\beta$ 满足 $(\mathrm{AR})_c$.
\end{enumerate}
则存在分次 $\operatorname{Gr}_{\mathfrak m}(A)$-模的同构 $\operatorname{Gr}_{\mathfrak m}(\operatorname{coker}\beta)\to\operatorname{Gr}_{\mathfrak m}(\operatorname{coker}\beta')$.
\end{lemma}
\begin{proof}
证明虽有些技术性, 却相当直接. 首先, 对 $L$ 和 $L'$ 取自由呈示, 可假设 $L=L'$. 可证明 $\beta'$ 满足 $(\mathrm{AR})_c$, 且第二个序列正合. 随后建立等式
\[
\mathfrak m^{n+1}N+\beta(M)\cap\mathfrak m^nN=\mathfrak m^{n+1}N+\beta'(M)\cap\mathfrak m^nN:
\]
用 $\beta$ 满足 $(\mathrm{AR})_c$ 证明包含关系 ``$\subseteq$'', 再用 $\beta'$ 满足 $(\mathrm{AR})_c$ 得到反向包含. 于是断言由描述 $\operatorname{Gr}_{\mathfrak m}(\operatorname{coker}\beta)_n=\mathfrak m^nN/(\mathfrak m^{n+1}N+\beta(M)\cap\mathfrak m^nN)$ 及 $\operatorname{Gr}_{\mathfrak m}(\operatorname{coker}\beta')_n$ 的类似描述推出. 详见 \cite[\S$3$]{CdJ02} 和 \cite[Tag 07VF]{SP}.
\end{proof}
\begin{proof}[Conrad--de Jong 逼近 (定理 \bookref{C.6.2}) 的证明]
因为 $\mathcal X$ 保持极限, 且 $R$ 是其有限生成 $\Bbbk$-子代数的余极限, 存在 $\Bbbk$ 上的有限型仿射概形 $V=\Spec B$ 及 $V$ 上的 $\mathcal X$ 的对象 $\gamma$, 配备 $2$-交换图
\[
\begin{tikzcd}
\Spec R\arrow[r]\arrow[rr,bend left=25,"\xi"]&V\arrow[r,"\gamma"']&\mathcal X.
\end{tikzcd}
\]
令 $v\in V$ 为极大理想 $\mathfrak m\subseteq R$ 的像. 必要时给环 $B$ 添加生成元, 可假设复合 $\widehat{\mathcal O}_{V,v}\to R\to R/\mathfrak m^2$ 满射. 由完备 Nakayama 引理 (\bookref{B.5.6}(3)), $\widehat{\mathcal O}_{V,v}\to R$ 满射. 现在的目标是在 $V$ 上同时逼近定义闭浸入 $\Spec R\hookrightarrow\Spec\widehat{\mathcal O}_{V,v}$ 的方程与关系, 以及对象 $\xi$. 为此, 选择 $\widehat{\mathcal O}_{V,v}$-模的消解
\begin{equation}\label{src:C.6.5}
\widehat{\mathcal O}_{V,v}^{\oplus r}\xrightarrow{\widehat\alpha}\widehat{\mathcal O}_{V,v}^{\oplus s}\xrightarrow{\widehat\beta}\widehat{\mathcal O}_{V,v}\to R\to0,
\end{equation}
并考虑函子
\[
\begin{aligned}
F:(\Sch/V)&\to\mathrm{Sets}\\
(T\to V)&\mapsto\{\text{复形 }\mathcal O_T^{\oplus r}\xrightarrow\alpha\mathcal O_T^{\oplus s}\xrightarrow\beta\mathcal O_T\}.
\end{aligned}
\]
不难验证此函子保持极限. \bookeqref{C.6.5} 的消解给出 $F(\widehat{\mathcal O}_{V,v})$ 的元素. 应用 Artin 逼近 (\bookref{B.5.18}), 得到 étale 态射 $(V'=\Spec B',v')\to(V,v)$ 及元素
\begin{equation}\label{src:C.6.6}
(B'^{\oplus r}\xrightarrow{\alpha'}B'^{\oplus s}\xrightarrow{\beta'}B')\in F(V'),
\end{equation}
使得 $\alpha',\beta'$ 模 $\mathfrak m^{N+1}$ 等于 $\widehat\alpha,\widehat\beta$.

令 $U=\Spec A\hookrightarrow\Spec B'=V'$ 为 $\operatorname{im}\beta'$ 定义的闭子概形, 令 $u=v'\in U$. 考虑复合
\[
\eta:U\hookrightarrow V'\to V\xrightarrow\gamma\mathcal X.
\]
因为 $R=\operatorname{coker}\widehat\beta$, $A=\operatorname{coker}\beta'$, 得到同构 $R/\mathfrak m^{N+1}\cong A/\mathfrak m_u^{N+1}$, 以及 $\xi|_{R/\mathfrak m^{N+1}}$ 与 $\eta|_{A/\mathfrak m_u^{N+1}}$ 的同构. 这给出了断言 (1)--(4).

为建立 (5), 需证明存在同构 $\mathfrak m^n/\mathfrak m^{n+1}\cong\mathfrak m_u^n/\mathfrak m_u^{n+1}$. 对 $n\leq N$, 同构 $R/\mathfrak m^{N+1}\cong A/\mathfrak m_u^{N+1}$ 保证这一点. 另一方面, $n\gg0$ 时, 可将其看作 Artin--Rees 引理 (\bookref{B.5.4}) 的推论. 为处理 $n$ 的中间范围, 需控制 Artin--Rees 引理中出现的常数. 首先注意, 应用 Artin 逼近之前可以增大 $N$, 确保 $\widehat\alpha,\widehat\beta$ 满足 $(\mathrm{AR})_N$. 因而可自由作此假设. 现在对 \bookeqref{C.6.5} 的正合复形 $\widehat{\mathcal O}_{V,v}^{\oplus r}\xrightarrow{\widehat\alpha}\widehat{\mathcal O}_{V,v}^{\oplus s}\xrightarrow{\widehat\beta}\widehat{\mathcal O}_{V,v}$, 以及将 \bookeqref{C.6.6} 限制到 $F(\widehat{\mathcal O}_{V,v})$ 得到的复形 $\widehat{\mathcal O}_{V,v}^{\oplus r}\xrightarrow{\widehat{\alpha'}}\widehat{\mathcal O}_{V,v}^{\oplus s}\xrightarrow{\widehat{\beta'}}\widehat{\mathcal O}_{V,v}$ 应用引理 \bookref{C.6.4}, 直接得到断言 (5). 另见 \cite{CdJ02} 和 \cite[Tag 07XB]{SP}.
\end{proof}
\begin{exercise}\label{src:C.6.7}
证明 Conrad--de Jong 逼近蕴含 Artin 逼近.
\end{exercise}
\subsection{Artin 代数化}\label{src:C.6.3-subsection}
Artin 代数化的结论比 Artin 逼近或 Conrad--de Jong 逼近更强, 因为不再需要逼近. 它保证存在 $\Bbbk$ 上的有限型带点仿射概形 $(\Spec A,u)$ 上的对象 $\eta$, 在所有阶数上均与给定的有效形式形变 $\xi$ 一致. 为确保这一点, 需要求 $\xi$ 在 $u$ 处\textit{万有}, 即限制 $\xi_n=\xi|_{A/\mathfrak m_u^{n+1}}$ 定义 $A$ 上的\textit{万有}形式形变 $\{\xi_n\}$ (定义 \bookref{C.4.2}).
\begin{theorem}[Artin 代数化]\label{src:C.6.8}
设 $\mathcal X$ 为 $\Sch/\Bbbk$ 上保持极限的预叠. 设 $(R,\mathfrak m)$ 为完备 Noether 局部 $\Bbbk$-代数, $\xi\in\mathcal X(R)$ 为有效的万有形式形变. 则存在
\begin{enumerate}[(1)]
\item $\Bbbk$ 上的有限型仿射概形 $\Spec A$ 及 $\Bbbk$-点 $u\in\Spec A$;
\item 对象 $\eta\in\mathcal X(A)$;
\item $\Bbbk$-代数的同构 $\alpha:R\xrightarrow\sim\widehat A_{\mathfrak m_u}$; 以及
\item 对 $n\geq0$ 的相容同构族 $\xi|_{R/\mathfrak m^{n+1}}\cong\eta|_{A/\mathfrak m_u^{n+1}}$ (在等同 $R/\mathfrak m^{n+1}\cong A/\mathfrak m_u^{n+1}$ 之下).
\end{enumerate}
\end{theorem}
注意, 我们并未断言 $\xi\cong\eta|_{\widehat A_{\mathfrak m_u}}$, 尽管这在某些情形成立, 例如 $\mathcal X$ 为代数叠时. 例 \bookref{C.7.17} 给出有效万有形式形变不唯一的例子, 此时存在两个在 étale 局部不同构的代数化.
\begin{remark}\label{src:C.6.9}
若已知 $R$ 为有限生成 $\Bbbk$-代数的完备化, 则本定理可视为 Artin 逼近的一个简单推论. 事实上, 取 $N=1$ 应用 Artin 逼近, 再用万有性得到相容的映射 $R\to A/\mathfrak m_u^{n+1}$, 从而得到模 $\mathfrak m^2$ 为同构的映射 $R\to\widehat A_{\mathfrak m_u}$. 因为 $R$ 与 $\widehat A_{\mathfrak m_u}$ 抽象同构, 由完备 Nakayama 引理 (\bookref{B.5.6}(3)), 同态 $R\to\widehat A_{\mathfrak m_u}$ 为同构, 断言随之得证. 一般情形的论证类似, 只是用 Conrad--de Jong 逼近代替 Artin 逼近.
\end{remark}
\begin{proof}
取 $N=1$ 应用 Conrad--de Jong 逼近 (\bookref{C.6.2}), 得到 $\Bbbk$ 上的有限型仿射概形 $\Spec A$ 及 $\Bbbk$-点 $u\in\Spec A$, 对象 $\eta\in\mathcal X(A)$, 同构 $\phi_2:\Spec A/\mathfrak m_u^2\to\Spec R/\mathfrak m^2$, 同构 $\alpha_2:\xi|_{R/\mathfrak m^2}\to\eta|_{A/\mathfrak m_u^2}$, 以及分次 $\Bbbk$-代数的同构 $\operatorname{Gr}_{\mathfrak m}(R)\cong\operatorname{Gr}_{\mathfrak m_u}(A)$. 我们断言, $\phi_2$ 和 $\alpha_2$ 可归纳延拓为相容的态射族 $\phi_n:\Spec A/\mathfrak m_u^{n+1}\to\Spec R$ 及同构 $\alpha_n:\xi|_{A/\mathfrak m_u^{n+1}}\to\eta|_{A/\mathfrak m_u^{n+1}}$. 事实上, 给定 $\phi_n$ 和 $\alpha_n$, $\xi$ 的万有性蕴含存在补全下述交换图的提升 $\phi_{n+1}$:
\[
\begin{tikzcd}
\Spec A/\mathfrak m_u^n\arrow[r,"\phi_n"]\arrow[d]&\Spec R\arrow[d,"\xi"]\\
\Spec A/\mathfrak m_u^{n+1}\arrow[r,"\eta|_{A/\mathfrak m_u^{n+1}}"']\arrow[ur,dashed,"\phi_{n+1}"]&\mathcal X.
\end{tikzcd}
\]
取极限得到同态 $\widehat\phi:R\to\widehat A_{\mathfrak m_u}$, 由完备 Nakayama 引理 (\bookref{B.5.6}(3)), 它满射. 另一方面, 对每个 $n$, $\Bbbk$-向量空间 $\mathfrak m^N/\mathfrak m^{N+1}$ 与 $\mathfrak m_u^N/\mathfrak m_u^{N+1}$ 维数相同. 这蕴含 $\widehat\phi$ 为同构. 另见 \cite[定理 $1.6$]{Art69b} 和 \cite[\S$4$]{CdJ02}.
\end{proof}
